\documentclass[10pt,a4paper]{amsart}
\usepackage[margin=19mm]{geometry}
\usepackage{amsmath,amssymb,mathtools}
\usepackage[scr=rsfs,cal=euler]{mathalfa}
\usepackage{xfrac}
\usepackage[unicode,hidelinks,bookmarks=false]{hyperref}
\newcommand{\ie}{i.e.\ }
\newcommand{\cf}{cf.\ }
\newcommand{\resp}{respectively\ }
\newcommand{\Subsection}{\subsection}
\providecommand{\nobreakdash}{\nobreak\hbox{-}}
\setlength{\emergencystretch}{2em}
\multlinegap0pt
\usepackage{bm}
\usepackage{bbm}
\usepackage{dsfont}
\usepackage{xspace}
\usepackage{cancel}
\usepackage{mathtools}
\usepackage{cleveref}
\usepackage{amsthm}
\makeatletter
\@ifundefined{claim}{\newtheorem{claim}{Claim}}{}
\@ifundefined{definition}{\newtheorem{definition}{Definition}}{}
\@ifundefined{lemma}{\newtheorem{lemma}{Lemma}}{}
\makeatother
\newcommand{\G}{\Gamma}
\newcommand{\TT}{TCC(a_1,a_2,\dots,a_s)}
\newcommand{\T}{TCC(a_1,a_2,\dots,a_{s-1})}
\title[Boxicity and Cubicity of Divisor Graphs and Power Graphs]{Boxicity and Cubicity of Divisor Graphs and Power Graphs}
\author{L. Sunil Chandran, Jinia Ghosh}
\date{Source: arXiv:2501.16233v2 (CC BY 4.0)}
\begin{document}
\maketitle

\noindent\emph{Source and licence.} Boxicity and Cubicity of Divisor Graphs and Power Graphs. Version arXiv:2501.16233v2 (CC BY 4.0), distributed under the Creative Commons Attribution 4.0 International licence, as recorded in the arXiv licence metadata for this exact version and shown on the abstract page https://arxiv.org/abs/2501.16233v2. Full rights notice: licenses/arxiv-2501.16233-CC-BY-4.0.txt.\par\smallskip

\noindent\textit{Editorial scope.} Counted unit: the Lemma whose source label is \texttt{inductive process} in the article (source0.tex lines 513--515, source label \texttt{inductive process}), delivered together with the article's own complete original proof of it (source0.tex lines 516--610, one \texttt{proof} environment of 95 lines). No mathematics is written, repaired or re-derived: statement and proof are the article's own text line for line. Required context retained uncounted: boxicity def 2 (lines 488--490). Every definition, display and label of those lines is retained unchanged. Omitted from this package and not redistributed: 1--487, 491--512, 611--795. Nothing inside a retained excerpt is omitted or altered: every retained body is delivered exactly as the article's lines are written, so no omission receipt of any kind accompanies this package. Citations: the retained excerpts carry no citation command at all, so no bibliography entry of the article is transcribed and no reference is redistributed. Reference adaptation: none was needed and none was made; every reference inside the delivered unit resolves inside it, so no unresolved reference or citation is produced. Printed slips kept literal: no printed word, symbol, hypothesis or number of the counted statement or of its proof was altered. Layout and class: the article's own class is \texttt{article} with packages including quoting, geometry, enumitem, amsthm, tikz, amsmath, hyperref, comment, lineno, array, cleveref, tocbibind, pifont, amssymb; the wrapper is the lane's pinned \texttt{amsart} editorial wrapper at 10pt on A4, which restates the theorem-like environments the retained text needs and adds the article's own macro definitions that the retained text needs (\texttt{\textbackslash G}, \texttt{\textbackslash T}, \texttt{\textbackslash TT}), with extra wrapper packages (bm, bbm, dsfont, xspace, cancel, mathtools, cleveref). The delivered numbering is the unit's own. Assets: no retained span contains a figure environment, a graphics-inclusion command, a PDF-page inclusion, a table or a TikZ picture; no image file is copied into this package, the package declares no asset dependency and no image file is redistributed with it. Licence: version arXiv:2501.16233v2 of this article is distributed under the Creative Commons Attribution 4.0 International licence, as recorded in the arXiv licence metadata for this exact version and shown on the abstract page https://arxiv.org/abs/2501.16233v2. A full rights notice is delivered at licenses/arxiv-2501.16233-CC-BY-4.0.txt. Characters: every retained excerpt is pure ASCII, so the delivered engine, XeLaTeX with its default Latin Modern text font, reports no missing character.
\begin{definition}\label{boxicity def 2}\textbf{(Alternate definition of boxicity, cubicity)} The \emph{boxicity (cubicity)} of a non-complete graph $\G$ is the minimum positive integer $k$ such that $\G = \cap_{i=1}^k I_i $, where $I_i$'s are (unit) interval graphs.
    
\end{definition}

\begin{lemma}\label{inductive process}
    $cub(\TT) \leq cub(\T)+a_s$, where $s \geq 2$.
\end{lemma}

\begin{proof} Let \( cub(\T) = k \). The proof involves constructing \( k + a_s \) unit interval graphs \( I_1, I_2, \dots, I_k,\) \(I_{k+1}, I_{k+2}, \dots, I_{k+a_s} \) such that \(\TT\) is the intersection of these unit interval graphs. That means our aim here is to show that each unit interval graph $I_i$, $1 \leq I \leq k+a_s$, is a supergraph of $\TT$, and every pair of non-adjacent vertices in $\TT$ is non--adjacent in one of the unit interval graphs. The first \( k \) unit interval graphs \( I_1, I_2, \dots, I_k \) are referred to as \textit{Type 1 unit interval graphs}, while the remaining \( a_s \) unit interval graphs \( I_{k+1}, I_{k+2}, \dots, I_{k+a_s} \) are called \textit{Type 2 unit interval graphs}. For this proof, a pair of non-adjacent vertices in $\TT$ is referred to as a \textit{non-edge} in \(\TT\). The classification of the \( k + a_s \) unit interval graphs corresponds to two distinct types of non-edges in \(\TT\), which are explained in detail in the remainder of the proof.

Clearly $V(\TT)=\{ (x,0),(x,1),\dots,(x,a_s) \ : \ x \in V(\T)\}$. By definition of $E(\TT)$, it is easy to see that two distinct vertices $(x,b)$ and $(y,c)$ are adjacent in $\TT$ if and only if $x \leq y$ and $a \leq b$.
    %\textcolor{blue}{but both the equalities do not hold simultaneously}.

     Let $\{(x,b),(y,c)\}$ be a non-edge in $\TT$. Note that, in this case, $(x,b)$ and $(y,c)$ are two distinct vertices in $\TT$, and moreover $x \neq y$. We categorize the non-edges of $\TT$ into the following two types:

\vspace{0.2cm}

    \textit{Type 1 non-edge}: Let $\{x,y\}$ be a non-edge in $\T$. That means $x$ and $y$ are two incomparable $(s-1)$-tuples. Then, irrespective of the value of $b$ and $c$, $(x,b)$ and $(y,c)$ are incomparable $s$-tuples and hence $\{(x,b),(y,c)\}$ is a non-edge in $\TT$. Note that when $s=2$, there is no non-edge of Type 1.

    \textit{Type 2 non-edge:} Let $\{x,y\}$ be an edge in $\T$. That means $x$ and $y$ are two distinct comparable $(s-1)$-tuples. Without loss of generality, assume that $x < y$. Then, there exists at least an index $i$ in $1 \leq i \leq s-1$ such that $x_i < y_i$. Let $b > c$. Since the $s$-th indices of $(x,b)$ and $(y,c)$ satisfy $b > c$, $\{(x,b),(y,c)\}$ is a non-edge in $\TT$.

    \vspace{0.2cm}

    We now define the Type 1 unit interval graphs \( I_1, I_2, \dots, I_k \). Since $cub(\T)=k$, \Cref{boxicity def 2} implies that $\T$ is the intersection of $k$ unit interval graphs, say $I'_1,I'_2,\dots, I'_k$. For each $1\leq i \leq k$, let $h_i$ be an interval representation of the unit interval graph $I'_i$. For each $1\leq i \leq k$, we define an interval representation $f_i$ of $I_i$ as follows: %The Type 1 unit interval graphs $I_i$, $1\leq i \leq k$ are defined as follows: for each $i$ in $1\leq i \leq k$, a mapping $f_i$ each $(x,b) \in V(\TT)$ to a closed unit interval $f_i((x,b))$ as follows: 

    
   \begin{center}
       $f_i((x,b)) = h_i(x)$, for all $x \in V(\T)$ and for all $0 \leq b \leq a_s$.
   \end{center} 

    \begin{claim}\label{type 1 interval graphs are supergraphs}
        For each $I_i$, $1 \leq i \leq k$, $E(\TT) \subseteq E(I_i)$.
    \end{claim}
    \begin{proof}
        % % In each $I_i$, for each fixed $x$, the set $\{ (x,a) \ | 0 \leq a \leq a_s\}$ forms a clique. 
        % As $I'_1,I'_2,\dots,I'_k$ is an interval graph representation of $\T$, so $E(\T)$ is a subset of $ E(I'_i)$, for all $1 \leq i \leq k$. Now in $I_i$, $f_i((x,a))=h_i(x)$, for all $x \in V(\T)$.
        
        Let $(x,b)$ and $(y,c)$ be adjacent in $\TT$. Then, $(x,b)$ and $(y,c)$ are comparable $s$-tuples. Without loss of generality, we assume that $x \leq y$ and $b \leq c$. This implies that either $x=y$ or $\{x,y\} \in E(\T)$. Note that in the both cases, $h_i(x) \cap h_i(y) \neq \emptyset$, for all $1 \leq i \leq k$. Moreover, $f_i((x,b))=h_i(x)$ and $f_i((y,c))=h_i(y)$. Hence, $f_i((x,b)) \cap f_i((y,c)) \neq \emptyset$ and therefore $(x,b)$ and $(y,c)$ are adjacent in each $I_i$.
       % Hence, $E(\TT) \subseteq E(I_i)$, for all $1 \leq i \leq k$.
    \end{proof}

   \begin{claim}\label{type 1 interval graphs and non-edge}
       % Each Type 1 non-edge of $\TT$ is a non-edge in one of the Type 1 unit interval graphs.
       For any Type 1 non-edge $\{(x,b),(y,c)\}$ of $\TT$,
    there exists some Type 1 unit interval graph $I_l$, such that $\{(x,b),(y,c)\} \notin E(I_l)$.
   \end{claim} 
    \begin{proof}
        Since $\{(x,b),(y,c)\}$ is a non-edge of Type 1, $\{x,y\}$ is a non-edge in $\T$. Hence there exists an unit interval graph $I'_l$ in $I'_1,I'_2,\dots,I'_k$ such that $\{x,y\} \notin E(I'_l)$, i.e., $h_l(x)\cap h_l(y)=\emptyset$. So, by our construction, $f_l((x,b))\cap f_l((y,c))=\emptyset$ and hence $\{(x,b),(y,c)\} \notin E(I_l)$.
    \end{proof}

From \Cref{type 1 interval graphs are supergraphs} and \Cref{type 1 interval graphs and non-edge}, one can see that the intersection $\bigcap_{i=1}^k I_i$ of the Type 1 unit interval graphs contains each edge of $\TT$, but no Type 1 non-edge of $\TT$.  % However, to determine the cubicity (or at least an upper bound of it) of $\TT$, additional conditions should be satisfied (see \Cref{boxicity def 2}). %For this purpose, in the rest of the proof, we use the term \textit{non-edge} in the following sense - $\{u,v\}$ is a non-edge in a graph $\G$ if $\{u,v\} \notin E(\G)$. 
% To compute an upper bound for $cub(\TT)$, we must identify a collection $\mathcal{C}$ of unit interval graphs such that any non-edge $\{(x,b),(y,c)\}$ in $\TT$ is a non-edge in at least one of the unit interval graphs in $\mathcal{C}$.
% \textcolor {red} {try to reduce the unnecessary use of "Now"}
However, each non-edge $\{(x,b),(y,c)\}$ of Type 2 is present in all Type 1 unit interval graphs of $\TT$. It is because in this case $\{x,y\} \in E(\T)$ and due to our construction, $f_i((x,b)) \cap f_i((y,c)) =h_i(x) \cap h_i(y) \neq \emptyset$, for all $1 \leq i \leq k$. To tackle this situation, we construct the Type 2 unit interval graphs $I_{k+1}, I_{k+2}, \dots, I_{k+a_s}$. Let $a_1+ a_2 + \dots + a_{s-1}=S$. For $x \in V(\T)$, let $|x|=\sum_{i=1}^{s-1}x_i$. Let $f_{k+a}$ be an interval representation of $I_{k+a}$, $1 \leq a \leq a_s$ defined as follows:

\begin{center}
    $f_{k+a}((x,a')) = [|x|,S+|x|]$, for all $0 \leq a' \leq a-1$,

     \hspace{1.4cm}        $= [|x|-S,|x|]$, for all $a \leq a' \leq a_s$.

\end{center}


%     Let $\{(x,b),(y,c)\}$ be a non-edge of Type (2). In that case, $\{x,y\}$ is a non-edge in $\T$, and hence there exists an unit interval graph $I'_i$ in $I'_1,I'_2,\dots,I'_k$ such that $\{x,y\} \notin E(I'_i)$. So, by our construction, $\{(x,b),(y,c)\} \notin E(I_i)$. So the $k$ unit interval graphs $I_1, I_2,\dots, I_k$ of $\TT$ together handle the non-edges of Type (2).  

     % \textcolor {red} {Are you sure that the reader will be able to understand what is meant by "handle" ? Have you made the reader aware what you are trying to do, in the beginning? Or are you writing only for the expert in Boxicity?}***************************
     
%     Whereas each non-edge $\{(x,b),(y,c)\}$ of Type (2) is present in $E(I_i)$ for all unit interval graphs $I_i$, $1 \leq i \leq k$, of $\TT$. 
%     \textcolor {red} {Are you using Whereas correctly? You are starting the sentence with it, without adding the part to which you are comparing the former part of the sentence. Just check this with grammarly. Try to intergrate grammarly with overleaf.}
%     It is because $\{x,y\} \in E(\T)$ and due to our construction, $f_i((x,b))=h_i(x),f_i((y,c))=h_i(y)$, for all $1 \leq i \leq k$. To tackle this situation, we construct $a_s$ many new interval graphs $I_{k+a}$, $1\leq a \leq a_s$. Let $a_1+ a_2 + \dots + a_{s-1}=S$. For $x \in V(\T)$, let $|x|=\sum_{i=1}^{s-1}x_i$. Let $f_{k+a}$ be an interval representation of $I_{k+a}$, $1 \leq a \leq a_s$ defined as follows:

%     $f_{k+a}((x,a')) = [|x|,S+|x|]$, for all $0 \leq a' \leq a-1 $

% \hspace{1.9cm} $= [|x|-S,|x|]$, for all $a \leq a' \leq a_s$.

Note that all the closed intervals in the interval representations defined above (for Type 2 unit interval graphs) have the same length $S$, and therefore, each closed interval can be scaled to a unit closed interval.

% Note that all the closed intervals in the interval representations above have the same length, $S$, and therefore, each closed interval can be scaled to a unit interval.


 % \textcolor {red} {My feeling is that  since Type 2 edges are easier to handle and they correspond to the first type of interval graphs, it is better to call them Type 1. Consider this.Also group  the interval graphs you produce into two types or two groups.. think about the reader..make it easy for them. }********************

\begin{claim}\label{Type 2 interval graphs are supergraphs}
    For each $I_{k+a}$, $1 \leq a \leq a_s$, $E(\TT) \subseteq E(I_{k+a})$.
\end{claim}
\begin{proof}
    Let $x$ be a vertex in $V(\T)$. It is easy to note that $|x| \in \{0,1,2,\dots,S\}$. For each $(x,a')$, where $0\leq a' \leq a-1$, the point $S \in f_{k+a} ((x,a'))$. Hence, the set $\{(x,a') \ : \ x \in V(\T), \ 0 \leq a' \leq a-1\}$ induces a clique in $I_{k+a}$. Similarly, for each $(x,a')$, where $a\leq a' \leq a_s$, the point $0 \in f_{k+a} ((x,a'))$. Hence, the set $\{(x,a') \ : \ x \in V(\T), \ a \leq a' \leq a_s\}$ induces a clique in $I_{k+a}$.

    Now, we consider the only remaining case. Let $(x,b)$ and $(y,c)$ are adjacent in $\TT$, where $x \leq y$, $0 \leq b \leq a-1$ and $a \leq c \leq a_s$. Note that, $f_{k+a}((x,b))=[|x|,S + |x|]$ and $f_{k+a}((y,b))=[|y|-S,|y|]$. Since $x \leq y$, we have $|x| \leq |y|$ and so $f_{k+a}((x,b)) \cap f_{k+a}((y,c))= [|x|,|y|]$. Hence, $\{(x,b),(y,c)\} \in E(I_{k+a})$.
\end{proof}

\begin{claim}\label{Type 2 interval graphs and non-edge}
    For any Type 2 non-edge $\{(x,b),(y,c)\}$ of $\TT$,
    there exists some Type 2 unit interval graph $I_{k+a}$ such that $\{(x,b),(y,c)\} \notin E(I_{k+a})$.
\end{claim}
\begin{proof}
    Since $\{(x,b),(y,c)\}$ is a non-edge of Type 2, without loss of generality, we can assume that $x >y$ and $b < c$. Now $b < c$ implies $b+1 \leq c$. By construction of the interval graph $I_{k+b+1}$, the vertices $(x,b)$ and $(y,c)$ are mapped as follows: $f_{k+b+1}((x,b))=[|x|,S+|x|]$ and $f_{k+b+1}((y,c))=[|y|-S,|y|]$. Since $x > y$ implies $|x| > |y|$, the intervals $[|x|,S+|x|]\cap [|y|-S,|y|] = \emptyset$ and hence, $(x,b)$ and $(y,c)$ are non-adjacent in $I_{k+b+1}$.
\end{proof}

  From \Cref{Type 2 interval graphs are supergraphs} and \Cref{Type 2 interval graphs and non-edge}, one can see that the intersection $\bigcap_{a=1}^{a_s}I_{k+a}$ of the Type 2 unit interval graphs contains each edge of $\TT$, but does not contain any Type 2 non-edge.
  
  From our discussion so far (\Cref{type 1 interval graphs are supergraphs} - \Cref{Type 2 interval graphs and non-edge}), we can conclude that $\TT=\bigcap_{i=1}^{k+a_s}I_i$. Hence, by \Cref{boxicity def 2}, $cub(\TT) \leq k+a_s$. 
\end{proof}

\end{document}
